\chapter*{Prefazione}
\addcontentsline{toc}{chapter}{Prefazione}

Questo libro nasce da una domanda semplice, posta guardando due fotografie affiancate: perché le rocce delle Ande sembrano parlare, coperte di segni lasciati da chi le ha abitate migliaia di anni fa, mentre le rocce delle Alpi venete sembrano tacere?

La domanda è ingannevole, e lo si scopre subito: non è chiaro se il silenzio delle Alpi venete sia reale o se sia il silenzio di chi non ha ancora ascoltato bene. Rispondere ha richiesto di allontanarsi più volte dalla domanda di partenza, per poi tornarci sopra con qualcosa in più.

Il metodo seguito in queste pagine è stato semplice da enunciare e non sempre facile da rispettare: ogni volta che una fonte offriva una spiegazione già pronta, si è cercato di risalire al lavoro originale su cui quella spiegazione si basava, invece di accontentarsi della sintesi. Non è un esercizio di pedanteria. Una sintesi, per sua natura, toglie i numeri, le eccezioni, i casi che non tornano: proprio le cose che permettono a chi legge di giudicare da sé quanto un'affermazione sia solida. Un manuale che dice ``i cacciatori preistorici sceglievano le selle per la loro posizione strategica'' dice qualcosa di vero, ma non permette di capire se quella scelta fosse una regola quasi assoluta o una tendenza fra tante, quanti siti la confermano, quanti la smentiscono, e con quali soglie di distanza, quota, visibilità. Solo tornando alla ricerca originale quei numeri riemergono, e con essi la possibilità di usarli per un problema nuovo, come quello posto in queste pagine per il Veneto.

Questo modo di procedere ha anche un'altra conseguenza, meno comoda ma più onesta: molte delle strade imboccate in questo lavoro si sono rivelate sbagliate, o comunque troppo semplici per il problema che si voleva risolvere. Si è pensato, in un primo momento, che il movimento degli animali potesse essere ricostruito con un semplice calcolo di costo energetico legato alla pendenza del terreno. Si è pensato che la vegetazione dell'epoca potesse essere un vincolo utile per restringere la ricerca. Si è pensato che i valichi di montagna fossero, di per sé, il segno più affidabile di un luogo frequentato in epoca preistorica. Ognuna di queste idee, messa alla prova dei dati reali e della letteratura scientifica, ha mostrato crepe che ne consigliavano l'abbandono o la correzione. Questo libro non nasconde quel percorso: lo racconta, perché il valore di un metodo si vede tanto nei vicoli ciechi che sa riconoscere quanto nelle strade che alla fine imbocca.

Il risultato non è un manuale definitivo, né un algoritmo pronto per trovare incisioni rupestri sulle Alpi venete. È piuttosto un tentativo di mostrare come si costruisce un'ipotesi di ricerca ragionevole quando i dati sono scarsi, le fonti divulgative sono generiche, e l'unico modo per procedere è tornare, un passo alla volta, alle domande più semplici: chi viveva qui, quando, come si muoveva, e perché.

Una nota pratica, per chi legge: le date in questo libro non sono espresse con le sigle convenzionali della letteratura archeologica (BP, cal BP e simili), la cui definizione tecnica \`e poco intuitiva per chi non lavora nel settore. Si \`e scelto invece di misurare il tempo rispetto all'anno zero del calendario comune, quello che tutti hanno incontrato a scuola. La conversione, con le sue approssimazioni dichiarate, \`e spiegata nel capitolo dedicato al tempo e al clima. Allo stesso modo, i termini tecnici della preistoria (nomi di culture, di fasi climatiche, di industrie litiche) sono stati evitati nel corpo del testo ogni volta che un'espressione più semplice poteva bastare; dove il termine tecnico resta necessario, compare fra parentesi alla prima occorrenza ed è raccolto, con la sua definizione, nel glossario finale.
